\section{Round 4: A strict asymptotic improvement beyond $1/\sqrt 8$}

\subsection{Round-4 objective}
We continue on path \textbf{(A)} (proof direction).
Round~3 showed that the entire \emph{variance-only} strategy is blocked at the constant
\(1/\sqrt 8\).
The most promising next step is therefore to add one genuinely new structural ingredient.
In this round we use a finite subset of the Fourier-sign constraints from the continuous
Swinnerton-Dyer--Haugland model, and show that these constraints already force a
\emph{strict} asymptotic improvement over the Round~2 bound.
The new conclusion is:
\[
 c\;>\;\frac{1}{\sqrt 8},
\]
where \(c=\lim_{N\to\infty} M(N)/N\) is the minimum overlap constant.
The gain is non-explicit, but rigorous.

\subsection{Round-3 foundation used}
We use the following Round~3 outputs as black boxes:
\begin{itemize}
\item Theorem~36.5: the universal lower bound \(M(N)\ge N/\sqrt{8-1/N^2}>N/\sqrt 8\).
\item Theorem~36.8: the constant \(1/\sqrt 8\) is a genuine barrier for any argument using only bounded variance of the difference law.
\item The qualitative lesson extracted in Round~3: any improvement beyond \(1/\sqrt 8\) must come from additional structure of the overlap function, not from variance alone.
\end{itemize}

\subsection{New insight / tool (Round-4)}
The new ingredient is a \emph{finite Fourier obstruction}.
Working in the continuous model for the minimum overlap constant, White records two key necessary properties of every overlap function \(M\):
\begin{enumerate}
\item \(\operatorname{Var}(M)\le 2/3\),
\item \(\int_{-2}^{2}\cos(\pi kx)M(x)\,dx\le 0\) for every integer \(k\ge 1\).
\end{enumerate}
Round~2 used only the first item.
We show here that adding only the four frequencies
\[
 k\in\{3,4,7,12\}
\]
already rules out every variance-extremizing density.
This yields a compactness proof that the optimal constant is not merely \(\ge 1/\sqrt 8\), but in fact \emph{strictly larger}.

\subsection{Attack plan (Round-4)}
The exact gap left after Round~3 is that the bound \(1/\sqrt 8\) might conceivably still be the true infimum of the continuous relaxation, even though the variance-only method cannot prove anything larger.
To close that gap we proceed as follows:
\begin{enumerate}
\item Pass to the continuous overlap model, which equals the asymptotic constant \(c\).
\item Relax to a larger compact class \(\mathcal C_4\) of densities satisfying only the necessary conditions:
nonnegativity, total mass \(1\), height at most \(1\), variance at most \(2/3\), and the four cosine-sign inequalities for \(k=3,4,7,12\).
\item Show that a minimizer of \(\|M\|_\infty\) over \(\mathcal C_4\) exists.
\item Show that if this minimum were \(1/\sqrt 8\), then the minimizer would have to be the uniform density on an interval of length \(2\sqrt 2\).
\item Prove that no such interval density can satisfy the four cosine-sign constraints.
\end{enumerate}
This overcomes the Round~3 obstruction: the variance extremizers are now excluded by an additional finite piece of structure.

\subsection{Work (Round-4)}
Let \(M(N)\) denote the discrete minimum-overlap quantity and let
\[
 c:=\lim_{N\to\infty} \frac{M(N)}{N}.
\]
We use the following external theorem.

\medskip
\noindent\textbf{External theorem (Swinnerton-Dyer--Haugland, as quoted by White).}
The constant \(c\) equals the infimum of \(\|M_f\|_\infty\) over all measurable
\(f:[-1,1]\to[0,1]\) with \(\int_{-1}^{1}f=1\), where
\[
 M_f(x):=\int_{-1}^{1} f(t)(1-f(x+t))\,dt,\qquad x\in[-2,2].
\]
Moreover every such \(M_f\) satisfies
\[
 0\le M_f(x)\le 1,\qquad \int_{-2}^{2} M_f(x)\,dx=1,
\]
\[
 \int_{-2}^{2}(x-E(M_f))^2M_f(x)\,dx\le \frac23,
\]
and
\[
 \int_{-2}^{2}\cos(\pi kx)M_f(x)\,dx\le 0\qquad (k\ge 1).
\]

Motivated by this, define \(\mathcal C_4\) to be the class of all measurable functions
\(M:[-2,2]\to[0,1]\) such that
\[
 \int_{-2}^{2} M(x)\,dx=1,
\qquad
 \int_{-2}^{2}(x-E(M))^2M(x)\,dx\le \frac23,
\]
and
\[
 \int_{-2}^{2}\cos(\pi kx)M(x)\,dx\le 0
 \qquad \text{for each }k\in\{3,4,7,12\},
\]
where
\[
 E(M):=\int_{-2}^{2} xM(x)\,dx.
\]
Since every genuine overlap function \(M_f\) lies in \(\mathcal C_4\), we have
\[
 c\ge \lambda_4,
 \qquad
 \lambda_4:=\inf_{M\in\mathcal C_4}\|M\|_\infty.
\]
We therefore seek a lower bound for \(\lambda_4\).

\medskip
\noindent\textbf{Lemma 36.10 (continuous variance bound with equality case).}
Let \(M\) be a probability density on \(\mathbb R\), i.e.
\(M\ge 0\), \(\int_{\mathbb R}M=1\), and assume \(\|M\|_\infty\le \Omega\).
Then
\[
 \operatorname{Var}(M)\ge \frac{1}{12\Omega^2}.
\]
Equality holds if and only if
\[
 M(x)=\Omega\,\mathbf 1_I(x)
\]
for some interval \(I\subset \mathbb R\) of length \(|I|=1/\Omega\).

\noindent\emph{Proof.}
Write \(\mu:=E(M)\) and let \(X\) be a random variable with density \(M\).
For every \(r\ge 0\),
\[
 \mathbb P(|X-\mu|\le r)=\int_{\mu-r}^{\mu+r}M(x)\,dx\le 2\Omega r.
\]
Hence
\[
 \mathbb P(|X-\mu|>r)\ge \max(0,1-2\Omega r).
\]
Using the tail-integral identity,
\[
 \operatorname{Var}(M)=\mathbb E[(X-\mu)^2]
 =\int_0^{\infty}2r\,\mathbb P(|X-\mu|>r)\,dr
 \ge \int_0^{1/(2\Omega)}2r(1-2\Omega r)\,dr
 =\frac{1}{12\Omega^2}.
\]
Now suppose equality holds.
Then equality must hold at every step above, so for every \(0\le r\le 1/(2\Omega)\),
\[
 \int_{\mu-r}^{\mu+r}M(x)\,dx=2\Omega r,
\]
and for every \(r\ge 1/(2\Omega)\),
\(
 \mathbb P(|X-\mu|>r)=0
\).
The second condition shows that \(M\) is supported on
\([\mu-1/(2\Omega),\mu+1/(2\Omega)]\).
The first condition implies that on each symmetric interval about \(\mu\), the average value of \(M\) is exactly \(\Omega\); since \(M\le \Omega\) a.e., this forces
\(M(x)=\Omega\) almost everywhere on
\([\mu-1/(2\Omega),\mu+1/(2\Omega)]\).
Because the total mass is \(1\), the interval length is exactly \(1/\Omega\).
Conversely, the uniform density on any interval of length \(1/\Omega\) has variance \(1/(12\Omega^2)\).
\hfill$\square$

\medskip
\noindent\textbf{Lemma 36.11 (interval extremizers violate a finite Fourier-sign set).}
For each real \(c\), define
\[
 J_c(x):=\frac{1}{2\sqrt 2}\,\mathbf 1_{[c-\sqrt 2,\,c+\sqrt 2]}(x).
\]
Assume that \([c-\sqrt 2,c+\sqrt 2]\subset[-2,2]\), equivalently
\(
 |c|\le 2-\sqrt 2
\).
Then there exists \(k\in\{3,4,7,12\}\) such that
\[
 \int_{-2}^{2}\cos(\pi kx)J_c(x)\,dx>0.
\]
In particular, no density of the form \(J_c\) belongs to \(\mathcal C_4\).

\noindent\emph{Proof.}
A direct computation gives, for every integer \(k\ge 1\),
\[
 \int_{-2}^{2}\cos(\pi kx)J_c(x)\,dx
 =\frac{1}{2\sqrt 2}\int_{c-\sqrt 2}^{c+\sqrt 2}\cos(\pi kx)\,dx
 =\frac{\sin(\pi k\sqrt 2)}{\pi k\sqrt 2}\cos(\pi kc).
\]
We now analyze signs.
Since
\[
 3\sqrt 2\in(4,9/2),\qquad 4\sqrt 2\in(11/2,6),\qquad 7\sqrt 2\in(19/2,10),\qquad 12\sqrt 2\in(16,17),
\]
we have
\[
 \sin(3\pi\sqrt 2)>0,
 \qquad
 \sin(4\pi\sqrt 2)<0,
 \qquad
 \sin(7\pi\sqrt 2)<0,
 \qquad
 \sin(12\pi\sqrt 2)>0.
\]
We partition the admissible range \(|c|\le 2-\sqrt 2\) into cases.

\smallskip
\emph{Case 1:} \(|c|\le 1/8\).
Then \(3\pi|c|\le 3\pi/8<\pi/2\), so \(\cos(3\pi c)>0\).
Since \(\sin(3\pi\sqrt 2)>0\), the \(k=3\) coefficient is positive.

\smallskip
\emph{Case 2:} \(1/8<|c|<3/8\).
Then \(4\pi|c|\in(\pi/2,3\pi/2)\), so \(\cos(4\pi c)<0\).
Since \(\sin(4\pi\sqrt 2)<0\), the \(k=4\) coefficient is positive.

\smallskip
\emph{Case 3:} \(3/8\le |c|<1/2\).
Then
\[
 7\pi|c|\in[21\pi/8,7\pi/2).
\]
Reducing modulo \(2\pi\), this lies in
\([5\pi/8,3\pi/2)\), where cosine is negative.
Hence \(\cos(7\pi c)<0\).
Since \(\sin(7\pi\sqrt 2)<0\), the \(k=7\) coefficient is positive.

\smallskip
\emph{Case 4:} \(|c|=1/2\).
Then
\[
 \cos(12\pi c)=\cos(6\pi)=1,
\]
and since \(\sin(12\pi\sqrt 2)>0\), the \(k=12\) coefficient is positive.

\smallskip
\emph{Case 5:} \(1/2<|c|\le 2-\sqrt 2\).
Because \(2-\sqrt 2<2/3\), we have
\[
 3\pi|c|\in(3\pi/2,2\pi),
\]
so \(\cos(3\pi c)>0\).
Again the \(k=3\) coefficient is positive.

All cases yield a positive cosine coefficient for some \(k\in\{3,4,7,12\}\).
\hfill$\square$

\medskip
\noindent\textbf{Theorem 36.12 (four Fourier inequalities already beat the variance barrier).}
The infimum \(\lambda_4=\inf_{M\in\mathcal C_4}\|M\|_\infty\) is attained, and satisfies
\[
 \lambda_4>\frac{1}{\sqrt 8}.
\]

\noindent\emph{Proof.}
Let
\[
 B:=\{M\in L^\infty([-2,2]): \|M\|_\infty\le 1\}.
\]
By Banach--Alaoglu, \(B\) is weak-* compact in \(L^\infty([-2,2])=(L^1([-2,2]))^*\).
Now set
\[
 K:=\{M\in B: 0\le M\le 1\ \text{a.e.}\}.
\]
We claim that \(K\) is weak-* closed in \(B\).
Indeed, if \(M_n\in K\) and \(M_n\rightharpoonup^* M\), then for every nonnegative \(\varphi\in L^1([-2,2])\),
\[
 \int_{-2}^{2}\varphi(x)M(x)\,dx=\lim_{n\to\infty}\int_{-2}^{2}\varphi(x)M_n(x)\,dx\ge 0,
\]
so \(M\ge 0\) a.e.; applying the same argument to \(1-M_n\) gives \(1-M\ge 0\) a.e.
Hence \(K\) is weak-* compact.
The class \(\mathcal C_4\) is weak-* closed inside \(K\), because the maps
\[
 M\mapsto \int_{-2}^{2}M(x)\,dx,
 \qquad
 M\mapsto \int_{-2}^{2}xM(x)\,dx,
 \qquad
 M\mapsto \int_{-2}^{2}x^2M(x)\,dx,
\]
and
\[
 M\mapsto \int_{-2}^{2}\cos(\pi kx)M(x)\,dx
 \qquad (k\in\{3,4,7,12\})
\]
are all weak-* continuous.
Hence \(\mathcal C_4\) is weak-* compact.
Since the \(L^\infty\)-norm is weak-* lower semicontinuous, there exists
\(M_*\in\mathcal C_4\) such that
\(
 \|M_*\|_\infty=\lambda_4
\).

Now Lemma~36.10 applied to \(M_*\) gives
\[
 \operatorname{Var}(M_*)\ge \frac{1}{12\lambda_4^2}.
\]
But \(M_*\in\mathcal C_4\), so \(\operatorname{Var}(M_*)\le 2/3\).
Therefore
\[
 \frac{1}{12\lambda_4^2}\le \frac23,
\qquad\text{hence}\qquad
 \lambda_4\ge \frac{1}{\sqrt 8}.
\]
Suppose for contradiction that \(\lambda_4=1/\sqrt 8\).
Then
\[
 \operatorname{Var}(M_*)\le \frac23 = \frac{1}{12\lambda_4^2},
\]
so equality holds in Lemma~36.10.
Hence
\[
 M_*(x)=\frac{1}{2\sqrt 2}\,\mathbf 1_{[c-\sqrt 2,\,c+\sqrt 2]}(x)
\]
for some real \(c\).
Because \(M_*\) is supported in \([-2,2]\), necessarily \(|c|\le 2-\sqrt 2\).
But Lemma~36.11 shows that such a density violates at least one of the four cosine-sign constraints defining \(\mathcal C_4\), contradiction.
Therefore \(\lambda_4>1/\sqrt 8\).
\hfill$\square$

\medskip
\noindent\textbf{Corollary 36.13 (strict asymptotic improvement).}
Let \(c\) be the minimum overlap constant.
Then there exists \(\eta>0\) such that
\[
 c\ge \frac{1}{\sqrt 8}+\eta.
\]
In particular,
\[
 c>\frac{1}{\sqrt 8}.
\]

\noindent\emph{Proof.}
Every genuine overlap function \(M_f\) belongs to \(\mathcal C_4\), so
\[
 c=\inf_f \|M_f\|_\infty\ge \inf_{M\in\mathcal C_4}\|M\|_\infty=\lambda_4.
\]
By Theorem~36.12, \(\lambda_4>1/\sqrt 8\).
Setting \(\eta:=\lambda_4-1/\sqrt 8\) gives the claim.
\hfill$\square$

\subsection{Adversarial verification}
\begin{itemize}
\item \emph{Use of Round~3.} The proof genuinely builds on Round~3's diagnosis that variance alone is insufficient. The new input is precisely the finite Fourier-sign set \(k\in\{3,4,7,12\}\).
\item \emph{Compactness step.} We minimize over the larger relaxed class \(\mathcal C_4\), not over the genuine overlap class. This is deliberate: a lower bound for the larger class is automatically a lower bound for the actual problem since every overlap function lies in \(\mathcal C_4\).
\item \emph{Lower semicontinuity.} The weak-* lower semicontinuity of the \(L^\infty\)-norm is standard. This is enough to ensure the infimum over the weak-* compact class \(\mathcal C_4\) is attained.
\item \emph{Equality case in Lemma~36.10.} The proof forces equality simultaneously in the tail bound and in the support truncation, which leaves only the uniform interval density. No other shape can attain the variance minimum under a height cap.
\item \emph{Coverage of all centers \(c\).} The case split in Lemma~36.11 covers the full admissible interval \(|c|\le 2-\sqrt 2\). Endpoints \(|c|=1/2\) are handled separately because the \(k=7\) coefficient vanishes there.
\item \emph{Scope of the conclusion.} The theorem proves a \emph{strict} asymptotic improvement, but it does not produce an explicit numerical value of \(\eta\). So it does not compete with the best known explicit lower bounds; it isolates the obstruction more sharply and proves that the Round~2/3 barrier is not the true asymptotic threshold.
\end{itemize}

\subsection{Final}
\textbf{UNRESOLVED (but strictly advanced).}
We have now proved a new theorem beyond Round~3:
there exists \(\eta>0\) such that the minimum overlap constant satisfies
\[
 c\ge \frac{1}{\sqrt 8}+\eta.
\]
Thus the Round~2 constant \(1/\sqrt 8\) is not merely a barrier of the variance-only method; it is \emph{not} the true asymptotic value even in the relaxed class obtained by adding only four Fourier-sign constraints.
The remaining gap is quantitative: we do not yet have an explicit lower bound for \(\eta\).

\subsection{Completion estimate}
\textbf{COMPLETION: 68\%}.

\subsection{References}
\begin{itemize}
\item Round~3 notes in \texttt{36\_solutions.tex}.
\item Ethan Patrick White, ``Erd\H{o}s' minimum overlap problem,'' arXiv:2201.05704.
\item J. K. Haugland, ``The minimum overlap problem revisited,'' arXiv:1609.08000 (for the Swinnerton-Dyer/Haugland continuous reformulation context).
\end{itemize}
